\documentclass[master.tex]{subfiles}

\begin{document}

% ======================== TITLE ================================
\ifSubfilesClassLoaded{%
  % Standalone title block
  \begin{center}
    {\LARGE\bfseries\color{isublue} Chapter 1: Introduction}\\[6pt]
    {\large Theory and Methods in Causal Inference}\\[4pt]
    {\normalsize Department of Statistics, Iowa State University}
  \end{center}
  \bigskip
  \hrule height 1.2pt \relax
  \smallskip
}{%
  % Master book: chapter heading
  \chapter{Introduction}
}

% ======================== LEARNING OBJECTIVES ==================
\begin{tcolorbox}[defstyle, title={Learning Objectives}]
By the end of this chapter, students should be able to:
\begin{enumerate}
  \item Articulate the distinction between an interventional distribution
        $f(y \mid \doop(T{=}t))$ and an observational conditional distribution
        $f(y \mid T{=}t)$, and explain why they differ whenever $T$ is
        endogenous.
  \item Describe the same causal question in all three languages --- SEM,
        DAG graph surgery, and potential outcomes --- and begin to translate
        between them.
  \item Derive the two densities explicitly in the Gaussian linear
        confounded model, and identify the endogeneity bias
        $\alpha\gamma\sigma_U^2/\sigma_T^2$ as the gap between them.
  \item State the two-step paradigm (identification then estimation) and
        locate each step in the causal inference pipeline.
  \item Distinguish the IV DAG from the mediation DAG and explain why
        conflating an instrument with a mediator leads to identification
        failure.
\end{enumerate}
\end{tcolorbox}

\medskip\noindent
This chapter is an orientation, not a complete treatment.  It introduces
four ideas --- the interventional/observational distinction, the causal
trinity, the Gaussian confounded model as a running example, and the
identification-versus-estimation paradigm --- that will be developed
rigorously in Chapters~2--4 and beyond.  After one reading, students
should have a map of the territory; the details will become clear as
the course proceeds.

% ==============================================================
\section{Motivation: Why Causal Inference Is Hard}
% ==============================================================

Causal inference studies when and how the observational distribution can be used to recover the distribution that would arise under interventions.

\subsection{The Central Question}

Causal inference is concerned with a deceptively simple question: what would
happen to $Y$ if we were to \emph{set} $T = t$?  This is fundamentally
different from the question that standard statistical procedures answer:
what is the distribution of $Y$ among units \emph{observed} to have $T = t$?

\begin{defn}{Interventional vs.\ Observational Distribution}
\begin{itemize}
  \item The \emph{interventional distribution} $f(y \mid \doop(T{=}t))$ is
        the distribution of $Y$ in a hypothetical world where $T$ has been
        \emph{externally set} to $t$, regardless of the usual
        data-generating mechanism for $T$.
  \item The \emph{observational conditional distribution} $f(y \mid T{=}t)$
        is the distribution of $Y$ among the subpopulation of units for
        whom $T = t$ was \emph{observed} to hold.
\end{itemize}
These two distributions coincide only when $T$ is randomized (exogenous).
In all other cases they differ, and the difference is \emph{endogeneity bias}.
\end{defn}


\begin{defn}{Exogeneity and Endogeneity} 
Suppose the outcome is generated by the structural equation
\[
Y = f_Y(T, X, U),
\]
where $T$ is the treatment, $X$ are observed covariates, and $U$ represents
unobserved common causes of $T$ and $Y$ (confounders).
\begin{itemize}
\item The treatment \(T\) is \emph{exogenous given \(X\)} if
\[
T \indep U \mid X .
\]
\item The treatment \(T\) is \emph{endogenous given \(X\)} if
\[
T \nindep U \mid X ,
\]
so that some unobserved variable \(U\) affects both \(T\) and \(Y\).
\end{itemize}
\end{defn}

\begin{example}{Drug Trial with Unmeasured Severity}
Does a new drug ($T=1$) reduce blood pressure ($Y$)?  Suppose sicker patients
are more likely to take the drug, so that unmeasured severity $U$ is a common
cause of $T$ and $Y$.  Then:
\begin{itemize}
  \item \emph{Observational study}: $\E[Y \mid T{=}1] > \E[Y \mid T{=}0]$.
        Drug users have higher blood pressure on average.  A na\"{i}ve analyst
        concludes the drug is harmful.
  \item \emph{Randomized trial}: $\E[Y \mid \doop(T{=}1)] < \E[Y \mid
        \doop(T{=}0)]$.  Forcing everyone to take the drug reduces blood
        pressure.  The drug is beneficial.
\end{itemize}
Same data.  Different questions.  Opposite answers.  The entire discrepancy
is caused by the back-door path $T \leftarrow U \to Y$, which the
observational quantity $\E[Y \mid T{=}1]$ cannot remove.
\end{example}

\begin{warn}{The Fundamental Problem of Causal Inference}
We observe $f(y \mid T{=}t)$ directly from data.  We want
$f(y \mid \doop(T{=}t))$.  The goal of \emph{identification theory} is to
find conditions under which the former can be used to recover the latter.
\end{warn}

% ==============================================================
\section{The Causal Trinity: Three Languages for One Idea}
% ==============================================================

The same causal question --- what happens when we set $T = t$? --- can be
expressed in three distinct but equivalent languages. Collectively, we call
these the \textbf{causal trinity}: the \emph{structural equation model} (SEM),
the \emph{directed acyclic graph} (DAG) with graph surgery, and the
\emph{potential outcomes} framework. Each member of the trinity illuminates a
different facet of the same underlying idea, and fluency in all three is
essential for modern causal reasoning. Figure~\ref{w1:fig:trinity} depicts their
relationships.

% ---------------------------------------------------------------
% Causal Trinity Figure
% ---------------------------------------------------------------
\begin{figure}[h]
\centering
\begin{tikzpicture}[
  lang/.style = {
    rectangle, rounded corners=6pt,
    draw, very thick,
    text width=3.4cm, align=center,
    minimum height=1.4cm,
    font=\bfseries\small
  },
  sem/.style   = {lang, draw=semcolor,  fill=semcolor!10,  text=semcolor},
  dag/.style   = {lang, draw=dagcolor,  fill=dagcolor!10,  text=dagcolor},
  po/.style    = {lang, draw=pocolor,   fill=pocolor!10,   text=pocolor},
  center/.style = {
    ellipse, draw=accent, fill=defbg, thick,
    text width=3.0cm, align=center,
    font=\small\itshape\color{accent}
  },
  trans/.style = {
    <->, thick, color=darkgrey,
    shorten >=4pt, shorten <=4pt
  },
  tlabel/.style = {font=\footnotesize\itshape, color=darkgrey,
                   fill=white, inner sep=2pt}
]

% Three vertices of the triangle
\node[po]  (PO)  at (90:4.2cm)   {Potential\\Outcomes\\$Y(t)$};
\node[sem] (SEM) at (210:4.2cm)  {Structural Equation\\Model (SEM)\\$Y = f_Y(T,U,\varepsilon)$};
\node[dag] (DAG) at (330:4.2cm)  {DAG \&\\Graph Surgery\\$\Gcal,\;\Gcal_T$};

% Central node
\node[center] (C) at (0,0) {What happens\\when we\\set $T = t$?};

% Edges: vertex to center (spokes)
\draw[-{Stealth[length=5pt]}, thick, color=pocolor!70, shorten >=4pt, shorten <=4pt]
  (PO)  -- (C);
\draw[-{Stealth[length=5pt]}, thick, color=semcolor!70, shorten >=4pt, shorten <=4pt]
  (SEM) -- (C);
\draw[-{Stealth[length=5pt]}, thick, color=dagcolor!70, shorten >=4pt, shorten <=4pt]
  (DAG) -- (C);

% Edges: translations between vertices
\draw[trans] (PO)  -- (SEM)
  node[tlabel, midway, left=2pt]
    {\parbox{2.6cm}{\centering consistency\\+ no interference}};

\draw[trans] (SEM) -- (DAG)
  node[tlabel, midway, below=4pt]
    {\parbox{2.6cm}{\centering equation\\$\leftrightarrow$ arrow}};

\draw[trans] (DAG) -- (PO)
  node[tlabel, midway, right=2pt]
    {\parbox{2.6cm}{\centering SUTVA\\+ graph structure}};

% Annotation: what each language excels at
\node[font=\scriptsize\color{pocolor}, text width=2.8cm, align=center]
  at ($(PO)+(0,1.2)$) {Defines estimands\\(ATE, ATT, LATE)};
\node[font=\scriptsize\color{semcolor}, text width=2.8cm, align=center]
  at ($(SEM)+(-1.6,-0.8)$) {Likelihood\\construction};
\node[font=\scriptsize\color{dagcolor}, text width=2.8cm, align=center]
  at ($(DAG)+(1.6,-0.8)$) {Identification\\via do-calculus};

\end{tikzpicture}
\caption{The \textbf{causal trinity}: three languages that express the same
causal question from different vantage points. Potential outcomes
define \emph{what} we want to estimate; the DAG with graph surgery specifies
\emph{how} to identify it from observed data; the SEM provides the
\emph{generative mechanism} for likelihood construction. Double-headed arrows
show the translation assumptions connecting adjacent languages. All three
converge on the same central question (ellipse).}
\label{w1:fig:trinity}
\end{figure}

\subsection{Language 1: Structural Equation Models}

A structural equation model (SEM) represents the causal data-generating
mechanism as a system of equations, one per variable, together with a joint
distribution over exogenous errors.

\begin{defn}{Structural Equation Model (SEM)}
A recursive SEM over variables $(U, T, Y)$ is a system
\[
  U \sim p(u), \qquad
  T = f_T(U, Z, \delta), \qquad
  Y = f_Y(T, U, \varepsilon),
\]
where $\delta$ and $\varepsilon$ are mutually independent exogenous
disturbances, $Z$ is an observed instrument, and $U$ is an unobserved
confounder.
\end{defn}

\begin{remark}[Error Independence in the NPSEM-IE]
\label{w1:rem:npsem-ie}
The assumption that the disturbances $\delta$ and $\varepsilon$ are
\emph{mutually independent} is a substantive restriction, not merely a
technical convenience.  Because $U$ is unobserved and enters both
$f_T$ and $f_Y$, error independence does not remove the dependence
between $T$ and $Y(t)$ induced by $U$; it restricts only the additional
channel of dependence running through the equation-specific errors
$(\delta,\varepsilon)$.  \citet{richardson2014ace} call SEMs with this
structure \emph{Nonparametric Structural Equation Models with Independent
Errors} (NPSEM-IEs).

These notes adopt the NPSEM-IE framework throughout for two practical
reasons: error independence ensures the DAG's Markov condition holds,
making the Markov factorization of the joint density immediate; and it
integrates cleanly with the graphical identification theory developed in
later chapters.  The identification results we derive --- the back-door
criterion, the do-calculus, the IV assumptions --- hold under weaker
conditions; the error-independence restriction is a modeling choice, not
an identification necessity.
\end{remark}

\medskip
\noindent\textbf{Observational regime.}
$T$ is determined by $f_T$.  Because $U$ appears in both $f_T$ and $f_Y$,
we have $\mathrm{Cov}(T,\varepsilon) \neq 0$ in general: this is
endogeneity.

\medskip
\noindent\textbf{Interventional regime.}
The do-operator replaces the equation for $T$ with a constant $t$, yielding
the \emph{mutilated system}:
\[
  U \sim p(u), \qquad T = t \text{ (fixed)}, \qquad Y = f_Y(t, U, \varepsilon).
\]
Now $T$ and $U$ are independent by construction: the back-door path through
$U$ has been severed.  The interventional density $f(y \mid \doop(T{=}t))$
is therefore the distribution of $f_Y(t, U, \varepsilon)$ when $U$ is drawn
from its marginal $p(u)$, not from the conditional $p(u \mid T{=}t)$.

\subsection{Language 2: Directed Acyclic Graphs}

A DAG makes the same surgery graphically visible.

\begin{defn}{Graph Surgery}
In the observational DAG, arrows into $T$ encode the mechanisms that
determine $T$ (including back-door paths through unobserved variables).
Under $\doop(T{=}t)$, all arrows into $T$ are \emph{deleted}.  The resulting
mutilated graph $\Gcal_T$ has no back-door paths from $T$ to $Y$:
the causal effect flows only along the direct path $T \to Y$.
\end{defn}

The observational and mutilated DAGs for the Gaussian confounded model are shown
side by side in Figure~\ref{w1:fig:surgery} in Section~\ref{w1:sec:iv}, where the structural
equations are introduced concretely.

\subsection{Language 3: Potential Outcomes}

\begin{defn}{Potential Outcome \citep{rubin1974estimating}}
$Y(t)$ is the outcome that \emph{would have been observed} had $T$ been set
to $t$, possibly contrary to fact.

\medskip\noindent
\textbf{SUTVA} (stable unit treatment value assumption): $Y_i = Y_i(T_i)$,
i.e., the observed outcome for unit $i$ equals the potential outcome at the
actual treatment received.  This rules out interference between units and
multiple versions of treatment.
\end{defn}

\noindent\textbf{Connection to the do-operator.}
Under standard structural assumptions (consistency + no interference), the
potential outcome $Y(t)$ and the interventional distribution are linked by
\[
  Y(t) \;\sim\; f(y \mid \doop(T{=}t)).
\]

\begin{remark}
Potential outcomes define what we want to estimate
(e.g., $\E[Y(1)] - \E[Y(0)]$), but they do not specify how to write an
observed-data likelihood for estimation.  The do-operator fills this gap by
providing an explicit link between the interventional density and the
observational distribution via identification rules (back-door, front-door,
do-calculus).
\end{remark}

\subsection{Roadmap: How the Trinity Organizes These Notes}

The three languages are not used equally throughout the course.  Each
Part is dominated by one language, with the others appearing as
supporting tools.  Understanding this organization in advance will help
you see why each language is introduced when it is.

\medskip
\noindent\textbf{Part~I --- Foundations: the language of DAGs.}
Part~I builds the graphical foundation.  The DAG and the do-operator
are the primary language here because they make the
interventional/observational distinction \emph{syntactically explicit}:
confounding is visible as an open back-door path, and identifying
assumptions are visible as graph criteria (back-door, front-door,
do-calculus).  A mistake that is easy to commit silently in
potential-outcomes notation becomes a type error in do-notation, and is
therefore caught before it propagates.  The SEM appears as the
generative backbone wherever closed-form calculations are needed,
particularly in the Gaussian confounded running example.

\medskip
\noindent\textbf{Part~II --- Identification: bridging DAGs and observed data.}
Part~II introduces the potential outcomes framework and shows how the
graphical identification conditions of Part~I translate into concrete
statistical objects.  Ignorability ($Y(t) \indep T \mid X$) is the
potential-outcomes counterpart of the back-door criterion; the
propensity score $\pi(X) = P(T{=}1 \mid X)$ is the object that
operationalizes it in observed data.  The IV assumptions (relevance,
exogeneity, exclusion) are first stated graphically in Part~I and then
re-expressed as moment conditions and compliance-type restrictions in
Part~II.  This Part is where the two languages are most explicitly
translated into each other.

\medskip
\noindent\textbf{Part~III --- Estimation: semiparametric and machine-learning methods.}
Once identification has been established in the language of either DAGs
or potential outcomes, Part~III turns to estimation.  The semiparametric
efficiency framework --- efficient influence functions, doubly robust
estimators, double machine learning --- is largely language-neutral: it
works with whatever moment condition the identification step produces.
Fluency in all three languages is expected by this point, as
Part~III draws freely on all of them.

\subsection{Historical Roots of the Causal Trinity}
\label{w1:subsec:history}

The three languages did not emerge from a single research program.  They
grew independently in different disciplines over more than a century, driven
by different scientific questions and different communities.  Understanding
this history explains why the three schools sometimes use incompatible
terminology for the same idea, and why unifying them is intellectually
non-trivial.

\medskip
\noindent\textbf{\color{semcolor}Structural Equation Models: genetics and econometrics.}
The oldest strand begins with the biologist Sewall Wright (1921), who
introduced \emph{path analysis} to decompose correlations among hereditary
traits into direct and indirect causal contributions --- arguably the first
use of a directed graph for causal reasoning.  The framework was
transplanted into economics by Haavelmo (1943), who argued that economic
relationships should be modeled as autonomous structural equations robust
to policy interventions --- an early statement of what Pearl would later
call ``graph surgery''.  The Cowles Commission (Koopmans, Klein, and others)
formalized simultaneous equation systems throughout the 1940s--50s, and
this SEM tradition became the standard language of causal inference in
econometrics.  Lucas's (1976) critique --- that reduced-form equations break
down when policy changes the data-generating process --- is essentially a
statement of the interventional/observational distinction, decades before
that phrase existed.  SEMs were independently adopted by sociologists and
psychologists (LISREL, factor analysis), making this the most
interdisciplinary of the three schools.

\medskip
\noindent\textbf{\color{pocolor}Potential Outcomes: experimental statistics and epidemiology.}
Neyman (1923) introduced the notation $Y_i(t)$ in the context of randomized
agricultural experiments, asking what a plot's yield \emph{would have been}
under an unobserved treatment.  Fisher (1935) placed randomization at the
foundation of causal inference via the randomization distribution.  The
crucial extension to \emph{observational} studies was made by Rubin (1974),
who formalized the assignment mechanism and stated SUTVA explicitly.
Holland's (1986) \textit{JASA} paper ``Statistics and Causal Inference''
brought the framework to mainstream statistics under the dictum ``no
causation without manipulation''.  Robins (1986) extended potential outcomes
to longitudinal treatments and time-varying confounding via $g$-computation.
Imbens and Angrist (1994) connected the framework to instrumental variables
and defined the LATE, bringing potential outcomes into empirical economics.
Today the framework is dominant in statistics, epidemiology, and much of
economics.

\medskip
\noindent\textbf{\color{dagcolor}DAGs and do-calculus: computer science and philosophy.}
Pearl developed Bayesian networks through the 1980s as a representation for
probabilistic reasoning in artificial intelligence.  The pivotal step came
in Pearl (1995), where the \emph{do-operator} was introduced: a formal
syntax for distinguishing interventional from observational distributions
within a graph-based algebra.  The 2000 monograph \textit{Causality}
synthesised DAGs, the do-calculus, identification theory, and the
connections to potential outcomes into a unified framework.  In parallel,
Spirtes, Glymour, and Scheines (1993) --- working from \emph{philosophy} at
Carnegie Mellon --- developed constraint-based algorithms (PC, FCI) for
learning causal structure from data.  The DAG framework was rapidly adopted
by epidemiology (Hern\'{a}n and Robins) and is now the primary language of
causal inference in computer science and AI.

\medskip
\noindent\textbf{Convergence.}
By the 2000s it became clear that all three frameworks are, under standard
assumptions, \emph{logically equivalent}: any result expressible in one can
be translated into the others.  The differences are in notation, native
strengths, and disciplinary culture --- not in the underlying mathematics.
The causal trinity is therefore not three competing theories but three
dialects of the same language, each optimised for a different task.

Figure~\ref{w1:fig:timeline} summarizes the key milestones chronologically
across the three schools.

% ---------------------------------------------------------------
% Timeline figure
% ---------------------------------------------------------------
\begin{figure}[ht]
\centering
\begin{tikzpicture}[
  every node/.style = {font=\scriptsize},
  milestone/.style  = {
    rectangle, rounded corners=3pt,
    draw, thick, text width=3.1cm, align=center,
    inner sep=3pt, minimum height=0.7cm
  },
  sem/.style = {milestone, draw=semcolor, fill=semcolor!10, text=semcolor!80!black},
  po/.style  = {milestone, draw=pocolor,  fill=pocolor!10,  text=pocolor!80!black},
  dag/.style = {milestone, draw=dagcolor, fill=dagcolor!10, text=dagcolor!80!black},
  tick/.style = {thick, color=darkgrey},
]

% ---- Main timeline axis ----
\draw[tick, -{Stealth[length=6pt]}] (0,0) -- (14.8,0);

% Year positions: x = (year - 1915) * 0.155
\foreach \yr/\xpos in {
  1920/0.77, 1930/2.32, 1940/3.87, 1950/5.42,
  1960/6.97, 1970/8.52, 1980/10.07, 1990/11.62, 2000/13.17}{
  \draw[tick] (\xpos, 0.12) -- (\xpos,-0.12)
    node[below=1pt, font=\tiny\color{darkgrey}] {\yr};
}

% ---- SEM milestones (top track, y = +1.85) ----
\node[sem] at ({(1921-1915)*0.155}, 1.85)
  {Wright (1921)\\path analysis};
\node[sem] at ({(1943-1915)*0.155}, 1.85)
  {Haavelmo (1943)\\structural eqs.};
\node[sem] at ({(1950-1915)*0.155}, 1.85)
  {Cowles Commission\\(1940s--50s)};
\node[sem] at ({(1976-1915)*0.155}, 1.85)
  {Lucas (1976)\\critique};

% ---- PO milestones (middle track, y = +0.65) ----
\node[po] at ({(1923-1915)*0.155}, 0.65)
  {Neyman (1923)\\$Y_i(t)$ notation};
\node[po] at ({(1935-1915)*0.155}, 0.65)
  {Fisher (1935)\\randomization};
\node[po] at ({(1974-1915)*0.155}, 0.65)
  {Rubin (1974)\\obs.\ studies};
\node[po] at ({(1986-1915)*0.155}, 0.65)
  {Holland (1986)\\JASA paper};
\node[po] at ({(1994-1915)*0.155}, 0.65)
  {Imbens--Angrist\\(1994) LATE};

% ---- DAG milestones (bottom track, y = -1.0) ----
\node[dag] at ({(1988-1915)*0.155}, -1.0)
  {Pearl (1988)\\Bayesian nets};
\node[dag] at ({(1993-1915)*0.155}, -1.0)
  {Spirtes et al.\\(1993) PC alg.};
\node[dag] at ({(1995-1915)*0.155}, -1.0)
  {Pearl (1995)\\do-operator};
\node[dag] at ({(2000-1915)*0.155}, -1.0)
  {Pearl (2000)\\\textit{Causality}};

% ---- Connector dots and stems ----
\foreach \yr in {1921,1943,1950,1976}{
  \fill[semcolor] ({(\yr-1915)*0.155}, 0) circle (2.2pt);
  \draw[semcolor!60, thin] ({(\yr-1915)*0.155}, 0.12)
    -- ({(\yr-1915)*0.155}, 1.22);
}
\foreach \yr in {1923,1935,1974,1986,1994}{
  \fill[pocolor] ({(\yr-1915)*0.155}, 0) circle (2.2pt);
  \draw[pocolor!60, thin] ({(\yr-1915)*0.155}, 0.12)
    -- ({(\yr-1915)*0.155}, 0.22);
}
\foreach \yr in {1988,1993,1995,2000}{
  \fill[dagcolor] ({(\yr-1915)*0.155}, 0) circle (2.2pt);
  \draw[dagcolor!60, thin] ({(\yr-1915)*0.155}, -0.12)
    -- ({(\yr-1915)*0.155}, -0.38);
}

% ---- Track labels ----
\node[font=\footnotesize\bfseries\color{semcolor}]
  at (-0.6, 1.85) {SEM};
\node[font=\footnotesize\bfseries\color{pocolor}]
  at (-0.6, 0.65) {PO};
\node[font=\footnotesize\bfseries\color{dagcolor}]
  at (-0.6,-1.00) {DAG};

% ---- Convergence annotation ----
\draw[thick, color=accent, decorate,
      decoration={brace, amplitude=5pt, mirror}]
  (12.5,-1.6) -- (14.5,-1.6)
  node[midway, below=6pt, font=\scriptsize\color{accent}, align=center]
    {convergence \&\\unification};

\end{tikzpicture}
\caption{Key milestones in the three schools of causal inference.
\textcolor{semcolor}{\textbf{SEM}} (top, blue) originates in genetics and
econometrics; \textcolor{pocolor}{\textbf{PO}} (middle, purple) in
experimental statistics; \textcolor{dagcolor}{\textbf{DAG}} (bottom, green)
in computer science and philosophy.  The three developed largely in isolation
until converging in the 2000s.}
\label{w1:fig:timeline}
\end{figure}

% ==============================================================
\section{The Core Distinction}
% ==============================================================

The single most important inequality of this course is:
\begin{equation}
  \underbrace{f\!\left(y \mid x,\, \doop(T{=}t)\right)}_{\text{interventional}}
  \;\neq\;
  \underbrace{f\!\left(y \mid x,\, T{=}t\right)}_{\text{observational}}
  \qquad \text{whenever } T \text{ is endogenous.}
  \label{w1:eq:core}
\end{equation}
The key assumption that connects the two densities is \emph{conditional
exchangeability}.

\begin{defn}{Conditional Exchangeability}
\label{w1:def:exch}
Given observed covariates $X$ and confounders $U$, the treatment assignment
is \emph{conditionally exchangeable with the intervention} if
\[
  Y(t) \indep T \mid X,\, U,
\]
or equivalently,
\[
  f(y \mid x,\, T{=}t,\, U{=}u) \;=\; f(y \mid x,\, \doop(T{=}t),\, U{=}u).
\]
This says that once we condition on all common causes $(X, U)$ of $T$ and
$Y$, observing $T{=}t$ and intervening to set $T{=}t$ produce the same
distribution of $Y$: there is no residual confounding given $(X, U)$.
\end{defn}

Under conditional exchangeability, the interventional density has a direct
representation in terms of observable quantities.

\begin{prop}{Back-Door Adjustment Formula}
\label{w1:prop:backdoor}
Under conditional exchangeability (Definition~\ref{w1:def:exch}),
\begin{equation}
  f(y \mid x,\, \doop(T{=}t))
  \;=\; \int f(y \mid x,\, T{=}t,\, U{=}u)\; p(u \mid x)\; du.
  \label{w1:eq:backdoor-formula}
\end{equation}
\end{prop}

\begin{proof}
\emph{Step~(i): weight under the intervention.}
Under $\doop(T{=}t)$, the arrow $U \to T$ is severed in the mutilated
graph $\Gcal_T$, so $U \indep T \mid X$ under the intervention.
The law of total probability therefore gives:
\[
  f(y \mid x,\, \doop(T{=}t))
  = \int f(y \mid x,\, \doop(T{=}t),\, U{=}u)\; p(u \mid x)\; du.
\]
\emph{Step~(ii): replacing the interventional kernel.}
By conditional exchangeability,
$f(y \mid x,\, \doop(T{=}t),\, U{=}u) = f(y \mid x,\, T{=}t,\, U{=}u)$.
Substituting completes the proof.
\end{proof}

\begin{remark}[Confounding Discrepancy]
\label{w1:rem:discrepancy}
What OLS estimates is the observational density $f(y \mid x, T{=}t)$,
not the interventional density.  Applying the law of total probability
to the observational density gives:
\[
  f(y \mid x,\, T{=}t)
  = \int f(y \mid x,\, T{=}t,\, U{=}u)\; p(u \mid x,\, T{=}t)\; du.
\]
The kernel $f(y \mid x, T{=}t, U{=}u)$ is the same as in
equation~\eqref{w1:eq:backdoor-formula}, but it is now averaged over
the \emph{selection-distorted} weight $p(u \mid x, T{=}t)$ rather than
the marginal $p(u \mid x)$.  Because $U$ and $T$ are dependent through the
back-door path, $p(u \mid x, T{=}t) \neq p(u \mid x)$, and the gap
\[
  f(y \mid x,\, T{=}t) - f(y \mid x,\, \doop(T{=}t))
  = \int f(y \mid x,\, T{=}t,\, U{=}u)
    \bigl[\, p(u \mid x,\, T{=}t) - p(u \mid x) \,\bigr]\, du
\]
is non-zero whenever $U$ and $T$ are dependent --- the endogeneity
condition.  This \emph{confounding discrepancy} motivates the two
scenarios below.
\end{remark}

\subsection{Two Scenarios: Observed vs.\ Unobserved Confounder}

Proposition~\ref{w1:prop:backdoor} expresses $f(y \mid x, \doop(T{=}t))$
as an integral over the observable kernel $f(y \mid x, T{=}t, U{=}u)$
weighted by $p(u \mid x)$.  Whether this formula can be \emph{used} from
data depends entirely on whether $U$ is observed.  The two scenarios lead
to fundamentally different identification strategies and motivate the
organization of this entire course.

\medskip
\noindent\textbf{Scenario 1: $U$ is observed (no unmeasured confounding).}
When every component of $U$ is recorded in the data, both
$f(y \mid x, T{=}t, U{=}u)$ and $p(u \mid x)$ can be estimated from the
observed sample $(Y, T, X, U)$.  The right-hand side of
equation~\eqref{w1:eq:backdoor-formula} is therefore a functional of the
observable distribution, and we call it the \emph{back-door adjustment
formula} (or standardization formula):
\begin{equation}
  f(y \mid x,\, \doop(T{=}t))
  = \int f(y \mid x,\, T{=}t,\, U{=}u)\; p(u \mid x)\; du.
  \label{w1:eq:backdoor-cont}
\end{equation}
The confounder is averaged over its population distribution $p(u \mid x)$,
not over the treatment-conditional distribution $p(u \mid x, T{=}t)$:
this is the adjustment that removes the confounding discrepancy.

A noteworthy special case of Scenario~1 is \emph{unconfoundedness}:
if $U \indep T \mid X$ already holds in the observational distribution,
then $p(u \mid x, T{=}t) = p(u \mid x)$ identically, and
formula~\eqref{w1:eq:backdoor-cont} reduces to the familiar observational
regression $\E[Y \mid X{=}x, T{=}t]$.  In other words, when conditioning
on $X$ alone suffices to remove confounding, no further adjustment is
needed: the observational distribution already equals the interventional
one.  Unconfoundedness holds by design in a randomized experiment, and
is the key assumption underlying regression adjustment and propensity
score methods studied in Chapter~6.

\medskip
\noindent\textbf{Scenario 2: $U$ is unobserved (unmeasured confounding).}
When $U$ is latent, neither $f(y \mid x, T{=}t, U{=}u)$ nor $p(u \mid x)$
can be estimated from data on $(Y, T, X)$ alone, and
formula~\eqref{w1:eq:backdoor-cont} is not usable.  Moreover, the
observational density $f(y \mid x, T{=}t)$ itself is a biased proxy for
the interventional density: as shown in Remark~\ref{w1:rem:discrepancy},
it averages the same kernel over the selection-distorted weight
$p(u \mid x, T{=}t)$ rather than $p(u \mid x)$, and the confounding
discrepancy $p(u \mid x, T{=}t) - p(u \mid x)$ cannot be estimated when
$U$ is unobserved.  This is the hard case, and it is the one that motivates
most of the advanced machinery in this course.

In Scenario~2, additional structure beyond $(Y, T, X)$ must be imposed to
recover $f(y \mid x, \doop(T{=}t))$.  The main identification strategies
studied in this course are \textbf{instrumental variables} (Chapter~7),
in which an observed variable $Z$ independent of $U$ provides
confounding-free variation in $T$; the \textbf{front-door criterion}
(Chapter~3), applicable when the effect of $T$ on $Y$ is fully mediated
by an unconfounded observed variable $M$; and \textbf{sensitivity
analysis}, which bounds conclusions by varying the assumed magnitude of
the confounding discrepancy $p(u \mid x, T{=}t) - p(u \mid x)$.

\noindent The two scenarios and their consequences are summarized in
Table~\ref{w1:tab:scenarios}.

\begin{table}[h]
\centering
\renewcommand{\arraystretch}{1.5}
\caption{The two scenarios arising from
Proposition~\ref{w1:prop:backdoor} and Remark~\ref{w1:rem:discrepancy}.}
\label{w1:tab:scenarios}
\begin{tabular}{@{}lll@{}}
\toprule
& \textbf{Scenario 1: $U$ observed} & \textbf{Scenario 2: $U$ unobserved} \\
\midrule
Confounding type
  & No \emph{unmeasured} confounding
  & Unmeasured confounding \\
Back-door formula usable?
  & Yes: $f(y\mid x,T{=}t,U{=}u)$ and $p(u\mid x)$ estimable
  & No: $U$ latent; neither quantity estimable \\
Identification route
  & Back-door adjustment / standardization
  & IV, front-door, \ldots \\
Key assumption
  & Unconfoundedness: $U \indep T \mid X$
  & Exclusion restriction, monotonicity, \ldots \\
Chapters in this course
  & Chapters~2--6 (back-door, IPW, DR)
  & Chapter~7 (IV) \\
\bottomrule
\end{tabular}
\end{table}

% ==============================================================
\section{The Gaussian Linear Confounded Model}
\label{w1:sec:iv}
% ==============================================================

\subsection{Setup}

The Gaussian linear confounded model is our primary running example
throughout this course.  It is simple enough to permit closed-form
calculations yet rich enough to illustrate every key concept.  We use it
here not to develop any particular identification strategy, but to make
the gap between observational and interventional reasoning mathematically
explicit; various identification strategies for this model appear in
later chapters.

\begin{example}{Gaussian Linear Confounded Model}
\label{w1:ex:ivmodel}
The structural equations are
\begin{equation}
  Y = \beta T + \gamma U + \varepsilon, \qquad T = \alpha U + \delta,
  \label{w1:eq:iv}
\end{equation}
where $U$ is an unobserved confounder and the three primitive variables
are mutually independent:
\[
  U \sim \mathcal{N}(0,\sigma_U^2), \qquad
  \varepsilon \sim \mathcal{N}(0,\sigma_\varepsilon^2), \qquad
  \delta \sim \mathcal{N}(0,\sigma_\delta^2).
\]
Here $\beta$ is the causal effect of $T$ on $Y$, $\gamma$ is the direct
effect of the confounder $U$ on $Y$, and $\alpha$ is the direct effect
of $U$ on $T$.  The error independence ($U \indep \varepsilon \indep
\delta$) is the NPSEM-IE condition of Remark~\ref{w1:rem:npsem-ie}.
\end{example}

The endogeneity of $T$ is immediate.  From the analyst's perspective, the
composite error in the $Y$ equation is $\gamma U + \varepsilon$, so:
\[
  \mathrm{Cov}(T,\; \gamma U + \varepsilon)
  = \gamma\,\mathrm{Cov}(\alpha U + \delta,\; U)
  = \gamma\alpha\sigma_U^2 \;\neq\; 0
  \quad (\text{when } \alpha \neq 0 \text{ and } \gamma \neq 0).
\]
Consequently, OLS regressing $Y$ on $T$ does not estimate $\beta$.

The causal structure is captured by the DAG pair in Figure~\ref{w1:fig:surgery}.
In the observational DAG, the back-door path $T \leftarrow U \to Y$ is open:
unobserved confounder $U$ causes both $T$ and $Y$, making $T$ and the
composite error $\gamma U + \varepsilon$ correlated.  Under the intervention
$\doop(T{=}t)$, all arrows into $T$ are deleted (shown faded), severing
the back-door path.

\begin{figure}[h]
\centering
\begin{tikzpicture}[node distance=1.8cm]
  %% --- Observational DAG ---
  \node[node]  (T1) at (0.5,0)    {$T$};
  \node[node]  (Y1) at (2.5,0)    {$Y$};
  \node[unode] (U1) at (1.5,-1.3) {$U$};
  \draw[edge]  (T1) -- (Y1);
  \draw[dedge] (U1) -- (T1);
  \draw[dedge] (U1) -- (Y1);
  \node[font=\small\itshape\color{darkgrey}] at (1.5,-2.2)
    {Observational DAG};
  %% --- Mutilated DAG ---
  \node[mnode] (T2) at (5.5,0)    {$t$};
  \node[node]  (Y2) at (7.5,0)    {$Y$};
  \node[unode] (U2) at (6.5,-1.3) {$U$};
  \draw[edge]  (T2) -- (Y2);
  \draw[dedge, opacity=0.25] (U2) -- (T2);   % severed
  \draw[dedge] (U2) -- (Y2);
  \node[font=\small\itshape\color{darkgrey}] at (6.5,-2.2)
    {Mutilated DAG: $\doop(T{=}t)$};
\end{tikzpicture}
\caption{Graph surgery on the Gaussian confounded model.  The faded arrow
in the mutilated DAG has been severed by the intervention $\doop(T{=}t)$.
The back-door path $T \leftarrow U \to Y$ is eliminated; only the causal
path $T \to Y$ remains active.  The treatment node becomes a fixed constant
(rectangle).}
\label{w1:fig:surgery}
\end{figure}

\subsection{Two Explicit Densities}

We now derive both densities in closed form.  Write $\sigma_T^2 =
\alpha^2\sigma_U^2 + \sigma_\delta^2$ for the marginal variance of $T$.

\medskip
\noindent\textbf{Interventional density.}
Set $T = t$ externally.  The back-door path through $U$ is severed, so
$U$ is independent of $T$ under the intervention and is drawn from its
marginal $\mathcal{N}(0,\sigma_U^2)$.  Then
$Y \mid \doop(T{=}t) = \beta t + \gamma U + \varepsilon$ with
$U \indep \varepsilon$.  Hence
\begin{equation}
  f\!\left(y \mid \doop(T{=}t)\right)
  = \mathcal{N}\!\left(\beta t,\;\;
    \gamma^2\sigma_U^2 + \sigma_\varepsilon^2\right).
  \label{w1:eq:int}
\end{equation}

\medskip
\noindent\textbf{Observational density.}
Observe $T = t$.  Since $Y$ and $T$ are both linear combinations of the
independent normals $(U, \varepsilon, \delta)$, the pair $(Y, T)$ is
jointly normal.  Their covariance is:
\[
  \mathrm{Cov}(Y, T)
  = \mathrm{Cov}(\beta T + \gamma U + \varepsilon,\; T)
  = \beta\sigma_T^2 + \gamma\alpha\sigma_U^2.
\]
Applying the conditional normal formula:
\begin{equation}
  f\!\left(y \mid T{=}t\right)
  = \mathcal{N}\!\left(
      \left(\beta + \frac{\gamma\alpha\sigma_U^2}{\sigma_T^2}\right)t,\;\;
      \frac{\gamma^2\sigma_U^2\sigma_\delta^2}{\sigma_T^2}
      + \sigma_\varepsilon^2
    \right).
  \label{w1:eq:obs}
\end{equation}

\subsection{The Endogeneity Gap}

Equations~\eqref{w1:eq:int} and~\eqref{w1:eq:obs} differ in both mean and variance:

\begin{center}
\renewcommand{\arraystretch}{1.4}
\begin{tabular}{@{}lll@{}}
\toprule
                   & \textbf{Interventional} & \textbf{Observational} \\
\midrule
Mean               & $\beta t$
                   & $\displaystyle\left(\beta +
                     \frac{\gamma\alpha\sigma_U^2}{\sigma_T^2}\right)t$ \\[6pt]
Variance           & $\gamma^2\sigma_U^2 + \sigma_\varepsilon^2$
                   & $\dfrac{\gamma^2\sigma_U^2\sigma_\delta^2}{\sigma_T^2}
                     + \sigma_\varepsilon^2$ \\[6pt]
Residual           & $(\gamma U{+}\varepsilon) \indep T$
                   & $(\gamma U{+}\varepsilon)$ correlated with $T$ \\
Identified by      & RCT (or IV, Chapter~7) & OLS \\
\midrule
\multicolumn{3}{@{}l@{}}{\small Equal when $\alpha = 0$ or $\gamma = 0$ (no confounding)} \\
\bottomrule
\end{tabular}
\end{center}

\medskip
The \emph{endogeneity bias} of OLS is $\gamma\alpha\sigma_U^2/\sigma_T^2$:
the coefficient of $T$ in the observational mean exceeds the causal
coefficient $\beta$ by this amount.  This is the omitted-variable bias
formula --- $U$ is omitted from the regression, it affects $Y$ via
$\gamma$ and $T$ via $\alpha$, and its contribution to the bias is scaled
by $\sigma_U^2/\sigma_T^2$.  Numerically, with $\alpha = \gamma = 1$ and
$\sigma_U^2 = \sigma_\delta^2 = 1$:
\[
  \text{bias} = \frac{1 \times 1 \times 1}{1^2 \times 1 + 1} = \frac{1}{2} = 0.5.
\]

\subsection{Lab: Simulating the Two Densities}

The analytical results of this section were verified by simulation using
$n = 10{,}000$ draws from the structural equations with parameters
$\beta = 2$, $\alpha = 1$, $\gamma = 1$,
$\sigma_U = \sigma_\varepsilon = \sigma_\delta = 1$.
The R code is distributed separately as \texttt{chapter1\_lab.R}.

\medskip
\noindent\textbf{Experiment 1: OLS bias.}
Data were generated from the structural equations with $\alpha=\gamma=1$
and all variances equal to 1.  OLS was applied to the observational pairs
$(Y_i, T_i)$.  By the omitted-variable bias formula, the OLS probability
limit is $\beta + \gamma\alpha\sigma_U^2/\sigma_T^2 = 2 + 0.5 = 2.5$.
The results confirm this:

\begin{center}
\renewcommand{\arraystretch}{1.3}
\begin{tabular}{@{}llll@{}}
\toprule
            & \textbf{Simulated} & \textbf{Theory} & \textbf{Formula} \\
\midrule
OLS slope   & 2.5147  & 2.5000  & $\beta + \gamma\alpha\sigma_U^2/\sigma_T^2$ \\
True $\beta$& ---     & 2.0000  & $\beta$ \\
Bias        & 0.5147  & 0.5000  & $\gamma\alpha\sigma_U^2/\sigma_T^2$ \\
\bottomrule
\end{tabular}
\end{center}

\noindent
OLS overestimates the causal effect by 25\%.

\medskip
\noindent\textbf{Experiment 2: Two-sample comparison at $t_0 = 1$.}
The observational sample consisted of all units with $|T_i - 1| < 0.15$;
the interventional sample was generated by setting $T = 1$ externally
and drawing fresh $U^*_i \sim \mathcal{N}(0,1)$ and
$\varepsilon^*_i \sim \mathcal{N}(0,1)$ independently.
Theoretical values are from equations~\eqref{w1:eq:int}
and~\eqref{w1:eq:obs} at $t_0 = 1$.

\begin{center}
\renewcommand{\arraystretch}{1.3}
\begin{tabular}{@{}lllll@{}}
\toprule
& \multicolumn{2}{c}{\textbf{Mean}} & \multicolumn{2}{c}{\textbf{SD}} \\
\cmidrule(lr){2-3}\cmidrule(lr){4-5}
& Simulated & Theory & Simulated & Theory \\
\midrule
$\E[Y \mid T{=}1]$         & 2.567 & 2.500
  & 1.240 & $\sqrt{1.5} = 1.225$ \\
$\E[Y \mid \doop(T{=}1)]$  & 1.999 & 2.000
  & 1.423 & $\sqrt{2} = 1.414$ \\
\bottomrule
\end{tabular}
\end{center}

\noindent
The variance reduction in the observational distribution reflects the
ratio $\sigma_{\mathrm{obs}}^2/\sigma_{\mathrm{int}}^2 = 1.5/2 = 0.75$:
conditioning on $T{=}1$ pins down the value of $\alpha U + \delta$,
which through the correlation with $\gamma U$ restricts the spread of $Y$.

\medskip
\noindent\textbf{Experiment 3: Density overlay for varying $\alpha$.}
Figure~\ref{w1:fig:density} shows the interventional density
$\mathcal{N}(\beta t_0,\, \gamma^2\sigma_U^2 + \sigma_\varepsilon^2)$
(solid) and the observational density \eqref{w1:eq:obs} (dashed),
both evaluated at $t_0 = 1$, for $\alpha \in \{0, 0.3, 0.7, 1.0\}$.
The interventional density does not depend on $\alpha$ and remains fixed;
as $\alpha$ increases, the observational curve shifts rightward
(growing bias) and narrows (reduced variance).

\begin{figure}[h]
\centering
\begin{tikzpicture}
\begin{scope}
  % Axis
  \draw[->] (-0.3,0) -- (7.5,0) node[right] {$y$};
  \draw[->] (0,-0.1) -- (0,2.8) node[above] {density};
  % x-axis ticks
  \foreach \x/\lab in {0/0, 1/1, 2/2, 3/3, 4/4, 5/5, 6/6, 7/7}
    \draw (\x,0.05) -- (\x,-0.05) node[below, font=\tiny] {\lab};
  % y-axis ticks
  \foreach \y/\lab in {0.5/0.1, 1.0/0.2, 1.5/0.3, 2.0/0.4, 2.5/0.5}
    \draw (0.05,\y) -- (-0.05,\y) node[left, font=\tiny] {\lab};

  % Interventional N(2,2) -- solid black; sqrt(2pi*2)=3.5449
  \draw[thick, black, domain=0:6, samples=80]
    plot (\x, {5*exp(-(\x-2)^2/4)/3.545});

  % alpha=0: same as interventional (grey dashes)
  \draw[thick, darkgrey, dashed, domain=0:6, samples=80]
    plot (\x, {5*exp(-(\x-2)^2/4)/3.545});

  % alpha=0.3: mean=2.2752, var=1.9174; sqrt(2pi*var)=3.4710
  \draw[thick, accent, dashed, domain=0:6, samples=80]
    plot (\x, {5*exp(-(\x-2.275)^2/3.835)/3.471});

  % alpha=0.7: mean=2.4698, var=1.6711; sqrt(2pi*var)=3.2404
  \draw[thick, accent!70!black, dashed, domain=0:6, samples=80]
    plot (\x, {5*exp(-(\x-2.470)^2/3.342)/3.240});

  % alpha=1.0: mean=2.5, var=1.5; sqrt(2pi*var)=3.0700
  \draw[thick, warnred, dashed, domain=0:6, samples=80]
    plot (\x, {5*exp(-(\x-2.5)^2/3.0)/3.070});

  % Legend
  \draw[thick, black]        (3.8,2.6) -- (4.4,2.6)
    node[right, font=\footnotesize] {Interventional $\mathcal{N}(2,2)$};
  \draw[thick, darkgrey, dashed]  (3.8,2.3) -- (4.4,2.3)
    node[right, font=\footnotesize] {Obs.\ $\alpha=0$};
  \draw[thick, accent, dashed]    (3.8,2.0) -- (4.4,2.0)
    node[right, font=\footnotesize] {Obs.\ $\alpha=0.3$};
  \draw[thick, accent!70!black, dashed] (3.8,1.7) -- (4.4,1.7)
    node[right, font=\footnotesize] {Obs.\ $\alpha=0.7$};
  \draw[thick, warnred, dashed]   (3.8,1.4) -- (4.4,1.4)
    node[right, font=\footnotesize] {Obs.\ $\alpha=1.0$};
\end{scope}
\end{tikzpicture}
\caption{Interventional density $\mathcal{N}(2,2)$ (solid black) versus
observational density for $\alpha \in \{0, 0.3, 0.7, 1.0\}$ (dashed),
evaluated at $t_0 = 1$.  Parameters: $\beta=2$, $\gamma=1$,
$\sigma_U=\sigma_\varepsilon=\sigma_\delta=1$.
As $\alpha$ increases, the observational curve shifts rightward (growing
bias) and narrows (reduced variance).  At $\alpha = 0$ there is no
confounding and the two densities coincide.
The interventional curve (solid) is invariant to $\alpha$.}
\label{w1:fig:density}
\end{figure}

% ---------------------------------------------------------------
% NEW SUBSECTION: IV DAG vs Mediation DAG
% ---------------------------------------------------------------
\subsection{The IV DAG Versus the Mediation DAG}
\label{w1:sec:iv_vs_med}

The IV DAG introduced above looks, at first glance, structurally similar to
the mediation DAG that will appear in Chapter~8.  Both involve a three-node
chain in which one variable sits between a ``cause'' and an ``effect''.  The
similarity is superficial; the conceptual difference is fundamental.
Confusing the two is one of the most common errors in applied causal
inference, and it is worth addressing head-on before the course proceeds.
This comparison is included only to prevent a common graphical confusion;
the full treatment of mediation analysis comes in Chapter~8.
Note that the instrument $Z$ introduced here for illustrative purposes
is not part of the running confounded model of this section; it will be
incorporated as an identification tool in Chapter~7.

Figure~\ref{w1:fig:iv_vs_med} places the two DAGs side by side.

\begin{figure}[H]
\centering
\begin{tikzpicture}[node distance=1.9cm]

  % ---- IV DAG (left panel) ----
  \begin{scope}[xshift=0cm]
    \node[node]  (Z1)  at (0,0)      {$Z$};
    \node[node]  (T1)  at (2.0,0)    {$T$};
    \node[node]  (Y1)  at (4.0,0)    {$Y$};
    \node[unode] (U1)  at (3.0,-1.4) {$U$};

    \draw[edge]  (Z1) -- (T1);
    \draw[edge]  (T1) -- (Y1);
    \draw[dedge] (U1) -- (T1);
    \draw[dedge] (U1) -- (Y1);

    % Exclusion restriction annotation
    \node[font=\scriptsize\color{warnred}, text width=2.8cm, align=center]
      at (2.0, 1.1)
      {\textit{exclusion restriction:\\$Z$ has no direct path to $Y$}};

    % Label
    \node[font=\small\bfseries\color{isublue}] at (2.0,-2.5)
      {(a)~Instrumental Variable DAG};
    \node[font=\footnotesize\color{darkgrey}, text width=4.6cm, align=center]
      at (2.0,-3.15)
      {$Z$ is \emph{upstream} of treatment,\\
       exogenous ($Z \indep U$), excluded from $Y$};
  \end{scope}

  % ---- Mediation DAG (right panel) ----
  \begin{scope}[xshift=7.2cm]
    \node[node]  (T2)  at (0,0)      {$T$};
    \node[node]  (M2)  at (2.0,0)    {$M$};
    \node[node]  (Y2)  at (4.0,0)    {$Y$};
    \node[unode] (U2)  at (2.0,-1.4) {$U$};

    \draw[edge]  (T2) -- (M2);
    \draw[edge]  (M2) -- (Y2);
    \draw[edge, bend left=30]  (T2) to (Y2);   % direct effect arc
    \draw[dedge] (U2) -- (T2);
    \draw[dedge] (U2) -- (Y2);

    % Indirect path annotation
    \node[font=\scriptsize\color{causalgreen}, text width=2.8cm, align=center]
      at (2.0, 1.1)
      {\textit{indirect path:\\$T \to M \to Y$}};

    % Label
    \node[font=\small\bfseries\color{isublue}] at (2.0,-2.5)
      {(b)~Mediation DAG};
    \node[font=\footnotesize\color{darkgrey}, text width=4.6cm, align=center]
      at (2.0,-3.15)
      {$M$ is \emph{downstream} of treatment,\\
       caused by $T$, lies on the causal path to $Y$};
  \end{scope}

\end{tikzpicture}
\caption{The IV DAG (a) and the mediation DAG (b) both involve a three-node
chain, but they are structurally and conceptually distinct.
In (a), the instrument $Z$ is pre-treatment and exogenous: it satisfies
$Z \indep U$ and has no direct path to $Y$ (exclusion restriction).
In (b), the mediator $M$ is post-treatment and endogenous: it is caused by
$T$ and transmits part of $T$'s effect to $Y$ (indirect path).
The curved arrow in (b) represents the direct effect $T \to Y$ bypassing $M$.
The unobserved confounder $U$ (dashed) opens a back-door path in both DAGs,
but the identification strategies it necessitates are entirely different.}
\label{w1:fig:iv_vs_med}
\end{figure}

Table~\ref{w1:tab:iv_vs_med} makes the comparison systematic.

\begin{table}[ht]
\centering
\renewcommand{\arraystretch}{1.5}
\caption{Instrumental variable versus mediator: a systematic comparison.}
\label{w1:tab:iv_vs_med}
\begin{tabular}{@{}p{3.8cm}p{5.0cm}p{5.0cm}@{}}
\toprule
& \textbf{Instrument} $Z$ & \textbf{Mediator} $M$ \\
\midrule
Position relative to $T$
  & Pre-treatment (upstream)
  & Post-treatment (downstream) \\
Causal relationship to $T$
  & $Z$ causes $T$ (first stage: $Z \to T$)
  & $T$ causes $M$ ($T \to M$) \\
Direct path to $Y$?
  & \textbf{No} (exclusion restriction)
  & \textbf{Yes} ($M \to Y$ is part of the total effect) \\
Independence from $U$?
  & \textbf{Yes} ($Z \indep U$, exogeneity)
  & Not required \\
Primary purpose
  & Identify $T \to Y$ in the presence of unmeasured confounding
  & Decompose $T \to Y$ into direct and indirect components \\
Key identifying assumptions
  & Relevance, exogeneity, exclusion
  & Cross-world independence;\\
  & & no $M$--$Y$ unmeasured confounding \\
Conditioning on it: effect
  & Controls endogeneity via first-stage
  & Can \emph{open} collider paths if $M$--$Y$ confounded \\
Course coverage
  & Chapter~7
  & Chapter~8 \\
\bottomrule
\end{tabular}
\end{table}

The two chains look similar but are not interchangeable:
\[
  \underbrace{Z \;\longrightarrow\; T \;\longrightarrow\; Y}_{%
    \text{IV chain: } Z \text{ is the identification device}}
  \qquad\text{versus}\qquad
  \underbrace{T \;\longrightarrow\; M \;\longrightarrow\; Y}_{%
    \text{Mediation chain: } M \text{ is on the causal path}}
\]
In the IV chain, $Z$ must satisfy $Z \indep U$ and the exclusion restriction
(no direct effect on $Y$); in the mediation chain, $M$ is caused by $T$ and
causes $Y$, and neither restriction applies.
Accordingly, a post-treatment variable should never be used as an instrument,
and a pre-treatment variable should never be interpreted as a mediator.

\begin{remark}
There is one important configuration where an observed post-treatment
variable serves a role \emph{analogous} to an instrument: the
\emph{front-door criterion} (Chapter~3).  When $T \to M \to Y$ and
$T \leftarrow U \to Y$ with $U$ unobserved, but $M$ is itself unconfounded,
the causal effect of $T$ on $Y$ can still be identified by combining two
back-door adjustments: first for the effect of $T$ on $M$, then for the
effect of $M$ on $Y$.  This is the one case where a mediator also plays an
identification role.  We will return to this in Chapter~3; for now, the IV
DAG and the mediation DAG should be understood as distinct structures
serving distinct purposes.
\end{remark}

% ==============================================================
\section{The Two-Step Paradigm}
\label{w1:sec:two-step}
% ==============================================================

Causal inference is a two-step discipline.  The two steps are logically
distinct and require different tools.

\subsection{Step 1: Identification}

\begin{defn}{Identification}
A causal quantity $P(y \mid \doop(T{=}t))$ is \emph{identified} from the
observational distribution $P$ if there exists a functional $\Psi$ such
that
\[
  P(y \mid \doop(T{=}t)) = \Psi(P),
\]
and $\Psi(P)$ depends only on the observable joint distribution.
\end{defn}

Identification is a purely mathematical question about the structure of the
causal model (DAG or SEM): given the assumed graph and assumptions, can the
interventional density be expressed as a function of the observed-data
distribution?  It does not depend on sample size. 

The main tools are:
\begin{itemize}
  \item the back-door criterion (introduced in Chapters~2--3, applied in Chapters~5--6),
  \item the front-door criterion (Chapter 3),
  \item the three rules of the do-calculus (Chapter 3),
  \item IV assumptions (Chapter 7).
\end{itemize}

\subsection{Step 2: Estimation}

Once $\Psi(P)$ has been identified, the statistical problem is to estimate
the functional $\Psi(P)$ from $n$ observations $O_1,\dots,O_n$ as
efficiently as possible.  The main tools are:
\begin{itemize}
  \item efficient influence functions (EIF, Part III),
  \item inverse probability weighting (IPW, Part III),
  \item doubly robust estimators (Part III),
  \item double machine learning (DML, Part III).
\end{itemize}

\subsection{A Worked Example: Flu Vaccination and Infection}
\label{w1:subsec:simpson}

The following example illustrates the two steps with a binary outcome and a
fully observed confounder, so no advanced machinery is required.

\medskip
\noindent\textbf{Scenario.}
A public health agency observes 1{,}000 individuals over one flu season.
Each person either received a flu vaccine ($T=1$) or did not ($T=0$), and
either became infected ($Y=1$) or did not ($Y=0$).  Age group $X$
(elderly $= E$, young $= Y$) is recorded for all individuals.
Elderly people are both more likely to be vaccinated (due to public health
campaigns) and more susceptible to infection (due to weaker immune systems),
so $X$ confounds the relationship between $T$ and $Y$.

The causal DAG is shown in Figure~\ref{w1:fig:vaccination}.  There are two
paths from $T$ to $Y$: the causal path $T \to Y$ (green), and the back-door
path $T \leftarrow X \to Y$ (red dashed).

\begin{figure}[h]
\centering
\begin{tikzpicture}[node distance=2.0cm]
  \node[node] (X) {$X$};
  \node[node, right of=X] (T) {$T$};
  \node[node, right of=T] (Y) {$Y$};
  \draw[cedge] (T) -- (Y) node[midway, above, font=\small\color{causalgreen}]
    {causal effect};
  \draw[bedge] (X) -- (T);
  \draw[bedge] (X) to[out=-30, in=-150] (Y);
  \node[font=\small\itshape\color{darkgrey}, below=0.3cm of X] {Age};
  \node[font=\small\itshape\color{darkgrey}, below=0.3cm of T] {Vaccine};
  \node[font=\small\itshape\color{darkgrey}, below=0.3cm of Y] {Infection};
  \node[font=\small\color{backdoorred}, below=1.0cm of T] {back-door path};
\end{tikzpicture}
\caption{Causal DAG for the vaccination example.  The green arrow is the
causal path $T \to Y$.  The red dashed arrows form the back-door path
$T \leftarrow X \to Y$: age $X$ causes both vaccination $T$ and infection
$Y$, creating a spurious association.  Because $X$ is observed, conditioning
on it blocks the back-door path.}
\label{w1:fig:vaccination}
\end{figure}

\medskip
\noindent\textbf{Observed data.}
The 1{,}000 individuals break down as follows (with row-conditional infection
rates in parentheses):

\begin{center}
\renewcommand{\arraystretch}{1.3}
\begin{tabular}{@{}lccccr@{}}
\toprule
& \multicolumn{2}{c}{\textbf{Vaccinated ($T=1$)}}
& \multicolumn{2}{c}{\textbf{Unvaccinated ($T=0$)}} & Total \\
\cmidrule(lr){2-3}\cmidrule(lr){4-5}
& $n$ & Infected & $n$ & Infected & \\
\midrule
Elderly ($X=E$) & 360 & 108\,(30\%) & 40  & 20\,(50\%) & 400 \\
Young ($X=Y$)   &  60 &   3\,(5\%)  & 540 & 54\,(10\%) & 600\\
\midrule
Total           & 420 & 111\,(26.4\%) & 580 & 74\,(12.8\%) & 1,000\\
\bottomrule
\end{tabular}
\end{center}

\medskip
\noindent\textbf{Simpson's paradox at a glance.}
The three infection-rate comparisons $\hat{P}(Y{=}1\mid T{=}1) - \hat{P}(Y{=}1\mid T{=}0)$
tell contradictory stories:

\begin{center}
\renewcommand{\arraystretch}{1.3}
\begin{tabular}{@{}lccc c@{}}
\toprule
\textbf{Stratum} & \textbf{Vaccinated} & \textbf{Unvaccinated} & \textbf{Difference} & \textbf{Conclusion} \\
\midrule
Elderly ($X=E$) & 30\% & 50\% & $-20$\,pp & vaccine helps \\
Young   ($X=Y$) &  5\% & 10\% & $-5$\,pp  & vaccine helps \\
\midrule
\textbf{Aggregate (raw)} & \textbf{26.4\%} & \textbf{12.8\%} & $\mathbf{+13.6}$\,\textbf{pp} & \textbf{vaccine harms??} \\
\bottomrule
\end{tabular}
\end{center}

\noindent
Within every stratum the vaccine reduces infection.  Yet in the aggregate
the vaccinated group has a \emph{higher} infection rate.  The reversal
occurs because the vaccinated group is disproportionately elderly (86\% of
vaccinated vs.\ 7\% of unvaccinated): a high-risk group is systematically
over-represented among the treated, making the raw comparison misleading.
The back-door formula corrects this by standardizing to the \emph{population}
age distribution rather than the treatment-conditional one.

\medskip
\noindent\textbf{Step 1: Identification.}
The only back-door path is $T \leftarrow X \to Y$.  The set $\{X\}$
satisfies the back-door criterion: it blocks this path (fork at $X$) and
contains no descendant of $T$.  By the back-door formula:
\begin{equation}
  P\!\left(Y{=}1 \mid \doop(T{=}t)\right)
  = \sum_{x \in \{E,Y\}}
    P(Y{=}1 \mid T{=}t, X{=}x)\,P(X{=}x).
  \label{w1:eq:backdoor}
\end{equation}
The causal quantity on the left has been expressed entirely as a functional
of the observed-data distribution on the right.

\medskip
\noindent\textbf{Step 2: Estimation.}
Plugging sample proportions into~\eqref{w1:eq:backdoor}, with
$\hat{P}(X{=}E) = 0.4$ and $\hat{P}(X{=}Y) = 0.6$:
\begin{align*}
  \hat{P}(Y{=}1 \mid \doop(T{=}1))
    &= 0.30 \times 0.4 + 0.05 \times 0.6 = 0.12 + 0.03 = 0.15, \\
  \hat{P}(Y{=}1 \mid \doop(T{=}0))
    &= 0.50 \times 0.4 + 0.10 \times 0.6 = 0.20 + 0.06 = 0.26.
\end{align*}
The estimated causal risk difference is
$\hat{\Delta}_{\mathrm{causal}} = 0.15 - 0.26 = -0.11$:
the vaccine causally reduces infection probability by 11 percentage points.

\medskip
\noindent\textbf{Why the na\"{i}ve comparison fails.}
The raw difference $0.264 - 0.128 = +0.136$ is precisely the aggregate
row of the Simpson table above: it mixes strata in treatment-dependent
proportions, importing the high elderly infection rate into the vaccinated
group.  The causal estimate $-0.11$ is correct because it re-weights each
stratum by $P(X{=}x)$, the population share, not by $P(X{=}x \mid T{=}t)$.

\begin{tcolorbox}[defstyle, title={The Two Steps, Concretely}]
\textbf{Step 1 (identification)} established --- from the DAG alone --- that
the back-door formula~\eqref{w1:eq:backdoor} expresses the causal quantity as
a functional of observable data.
\textbf{Step 2 (estimation)} replaced the population probabilities in that
formula with sample proportions.

These two steps are logically separate: Step 1 is a mathematical argument
about the causal model that does not depend on sample size; Step 2 is a
statistical argument about finite-sample estimation that does not depend on
the causal model.  A researcher who skips Step 1 and reports the raw
difference conflates the two and draws the wrong conclusion.
\end{tcolorbox}

% ==============================================================
\section{Summary}
% ==============================================================

\begin{enumerate}
  \item \textbf{Interventional vs.\ observational distribution.}
        Causal inference answers questions about interventions
        $f(y \mid \doop(T{=}t))$, not about observations $f(y \mid T{=}t)$.
        These two distributions coincide only under exogeneity (i.e.,
        $\alpha = 0$ or $\gamma = 0$ in the Gaussian confounded model); whenever $T$ is endogenous
        they differ, and the difference is endogeneity bias.

  \item \textbf{The causal trinity.}
        The same causal question can be expressed in three languages ---
        SEM (equation surgery), DAG (graph surgery), and potential outcomes
        ($Y(t)$).  The do-operator makes the interventional/observational
        distinction syntactically explicit; the other two languages leave
        it implicit.  Fluency in all three is essential for modern causal
        reasoning.

  \item \textbf{Endogeneity bias in the Gaussian confounded model.}
        In the Gaussian linear confounded model, the coefficient of $T$ in
        the observational mean $\E[Y \mid T{=}t]$ exceeds the causal
        coefficient $\beta$ by $\gamma\alpha\sigma_U^2/\sigma_T^2$, where
        $\sigma_T^2 = \alpha^2\sigma_U^2 + \sigma_\delta^2$.  This is the
        omitted-variable bias: $U$ drives both $T$ (via $\alpha$) and $Y$
        (via $\gamma$), and its omission from the regression inflates the
        OLS estimand.  The observational variance
        $\gamma^2\sigma_U^2\sigma_\delta^2/\sigma_T^2 + \sigma_\varepsilon^2$
        is smaller than the interventional variance
        $\gamma^2\sigma_U^2 + \sigma_\varepsilon^2$.

  \item \textbf{Two-step paradigm.}
        Causal inference is a two-step discipline: \emph{identification}
        (expressing $\Psi(P)$ as a functional of observable data, a
        mathematical question about the causal model) followed by
        \emph{estimation} (constructing $\hat{\Psi}_n$ efficiently from $n$
        observations, a statistical question about finite samples).  The two
        steps are logically independent.

  \item \textbf{Instrument vs.\ mediator.}
        The IV DAG ($Z \to T \to Y$ with $Z \indep U$) and the mediation
        DAG ($T \to M \to Y$) are structurally distinct.  The instrument is
        upstream of the treatment, exogenous, and excluded from the outcome
        equation; the mediator is downstream of the treatment, caused by
        $T$, and on the causal path.  Conflating the two leads to violations
        of the exclusion restriction and identification failure.
\end{enumerate}

\medskip
\noindent\textit{Causal inference is the study of when the observational
distribution contains enough information to recover the interventional
distribution.}  Everything in this course --- the back-door criterion,
do-calculus, instrumental variables, doubly robust estimation --- is an
answer to that question in a particular setting.


\section*{Problems}
\addcontentsline{toc}{section}{Problems}
% ==============================================================

\begin{enumerate}[label=\textbf{\arabic*.}, leftmargin=*, itemsep=10pt]

  \item \textbf{Do-notation fundamentals.}
  For each statement below, decide whether it refers to an interventional
  or an observational distribution, and rewrite it unambiguously using
  do-notation where appropriate.
  \begin{enumerate}[label=(\alph*)]
    \item ``The probability that a patient recovers given that they took the drug.''
    \item ``The probability that a patient would recover if we prescribed
          the drug to everyone.''
    \item ``Among students who attended tutoring, the average exam score
          was 85.''
    \item ``If we enrolled all students in tutoring, the average exam score
          would be 85.''
  \end{enumerate}

  \item \textbf{Deriving the observational density.}
  In the Gaussian confounded model (Example~\ref{w1:ex:ivmodel}), verify
  equation~\eqref{w1:eq:obs} by carrying out the following steps.
  Let $\sigma_T^2 = \alpha^2\sigma_U^2 + \sigma_\delta^2$.
  \begin{enumerate}[label=(\alph*)]
    \item Show that $(Y, T)$ is jointly normal by writing both as linear
          combinations of the independent normals $(U, \varepsilon, \delta)$.
    \item Compute $\mathrm{Cov}(Y, T)$ and $\mathrm{Var}(T)$.
    \item Apply the conditional normal formula to derive
          $\E[Y \mid T{=}t]$ and $\mathrm{Var}[Y \mid T{=}t]$,
          and hence verify equation~\eqref{w1:eq:obs}.
    \item Show that the OLS probability limit is
          $\beta + \gamma\alpha\sigma_U^2/\sigma_T^2$, and interpret
          this as an omitted-variable bias formula.
  \end{enumerate}

  \item \textbf{The causal trinity.}
  Consider the following verbal causal claim: ``Aspirin ($T$) reduces the
  risk of heart attack ($Y$) because it inhibits platelet aggregation ($M$),
  but patients with pre-existing cardiovascular disease ($U$, unobserved)
  are both more likely to take aspirin and more likely to have a heart
  attack.''
  \begin{enumerate}[label=(\alph*)]
    \item Draw the causal DAG implied by this description.
    \item Write down the recursive SEM with appropriate structural functions
          $f_T$, $f_M$, $f_Y$.
    \item State the SUTVA assumption and discuss whether it is plausible in
          this setting.
    \item Explain why $\E[Y \mid T{=}1] - \E[Y \mid T{=}0]$ does not identify
          the causal effect $\E[Y \mid \doop(T{=}1)] - \E[Y \mid \doop(T{=}0)]$
          in this DAG.
  \end{enumerate}

  \item \textbf{IV versus mediation.}
  For each scenario below, determine whether the intermediate variable
  described is an \emph{instrument} or a \emph{mediator}, and state which
  identifying assumption it satisfies or violates.
  \begin{enumerate}[label=(\alph*)]
    \item A researcher uses distance to the nearest college ($Z$) to study
          the effect of college attendance ($T$) on earnings ($Y$).
    \item A nutritionist studies how a high-fiber diet ($T$) reduces
          cholesterol ($M$), which in turn reduces heart disease risk ($Y$).
    \item A political scientist uses rainfall on election day ($Z$) to
          estimate the effect of voter turnout ($T$) on electoral outcomes
          ($Y$).
    \item A pharmacologist studies whether a drug ($T$) reduces blood
          pressure ($Y$) by inhibiting a specific enzyme ($M$), and
          proposes to use $M$ as an instrument for $T$.  Evaluate this
          proposal.
  \end{enumerate}

\end{enumerate}

% ==============================================================
\newpage
\ifSubfilesClassLoaded{
  \bibliographystyle{plainnat}
  \bibliography{causal_inference}
}{}
\end{document}